\section{Deep Learning for Radar Signal Analysis}\label{sec:related:dl_radar}

Learned models replace histogram rules with a parametric map from \gls{pdw} sequences or time-frequency images to labels.
Most published systems are supervised classifiers: they assume a global catalogue of emitter types, or of intra-pulse modulations, and they train a \gls{cnn} or \gls{rnn} with cross-entropy \parencite{qu2026radardeinterleavingreview}.
Attribute-specific recurrent networks on \gls{pdw} streams are a representative sequence model in this family \parencite{notaro2019emitter}.
They capture temporal dependence that a bag of pulse parameters misses, but the output is still a class index in a fixed label set.

A closer relative is supervised contrastive training for emitter classification.
\textcite{jonsson2023supcon} trains an encoder so that samples with the same emitter label lie close in embedding space, using synthetic scenarios from the same simulator lineage as \cref{sec:method:data}.
The loss belongs to the family used in \cref{sec:method:loss}, but the task is still closed-set identification: emitter identifiers are treated as globally comparable classes, and inference is classification rather than clustering.
Recurrent ensembles for emitter classification under dataset drift make the same closed-set assumption and treat out-of-distribution detection as a separate module \parencite{coleman2023rwrdrift}.

Deinterleaving as clustering, with an unknown number of emitters, is a different problem.
\textcite{gunn2025deinterleaving} train a transformer encoder with a triplet loss so that, within one interleaved stream, pulses from the same emitter are close and pulses from different emitters are far.
At inference they cluster the embeddings with hierarchical density clustering (HDBSCAN) \parencite{campello2015hierarchical,mcinnes2017hdbscan}.
They do not assume a known emitter count, they evaluate with \gls{ami}, \gls{ari}, and \gls{vmeasure}, and they omit index-based positional encodings because \gls{toa} already orders the pulses.
That pipeline is the closest published counterpart to \cref{sec:method:arch}.

Three differences remain.
First, Gunn et al.\ produce a single embedding and a single partition---emitter identity---and treat operating mode as a later analysis that becomes easier after deinterleaving.
Second, they train with a triplet hinge, which uses one positive and one negative per anchor; the method in \cref{sec:method:loss} uses the supervised contrastive kernel of \textcite{khosla2020supcon}, which averages over all positives of an anchor in the mini-batch.
Third, their transformer uses vanilla dot-product attention.
This thesis additionally evaluates a pairwise physical bias and a rotary encoding of time and azimuth (\cref{sec:method:attention}).

Deep models for radar pulses already exist.
What they still omit is the joint problem of this thesis: two embeddings whose positive sets have different scopes---recording-local emitters and globally shared modes---both decoded by clustering rather than by a catalogue softmax.
